\chapter{How the results are measured}
\label{ch:measuring}

Every number in Tracks~A--C comes from one of the measures in this chapter. Learn them here once, and the results chapters will read easily.

\section{Scoring a classification}

\subsection{The confusion matrix}

A \textbf{confusion matrix} is a table with the gold classes as rows and the predicted classes as columns. Correct answers lie on the diagonal. A model that gives the same answer to everything fills a single column. Confusion matrices are the quickest way to see \emph{how} a model is right or wrong, so they appear in every results chapter.

\subsection{Accuracy, precision, recall and F1}

For class $c$, let $TP_c$ be the patches of class $c$ predicted as $c$ (true positives), $FP_c$ the patches of other classes predicted as $c$ (false positives) and $FN_c$ the patches of class $c$ predicted as something else (false negatives). Then
\begin{equation}
\mathrm{Acc} = \frac{\#\text{correct}}{N}, \quad
P_c = \frac{TP_c}{TP_c+FP_c}, \quad
R_c = \frac{TP_c}{TP_c+FN_c}, \quad
F1_c = \frac{2P_cR_c}{P_c+R_c},
\label{eq:metrics}
\end{equation}
and
\begin{equation}
\text{macro-F1} = \frac{1}{3}\sum_{c} F1_c .
\label{eq:macrof1}
\end{equation}

\begin{description}
  \item[Accuracy] The share of test patches answered correctly. Easy to read, but it can hide a class the model never gets right.
  \item[Macro-F1] The average of the three per-class F1 scores; the roadmap's primary metric. Every class counts equally, so a model that ignores one class is punished even if its accuracy looks fine.
  \item[Majority baseline] Always answering the most common test class: 34.5\,\% accuracy, macro-F1 0.171.
\end{description}

\begin{plainwords}
With three nearly balanced classes, guessing gives about 33\,\%. On 55 patches, any accuracy between about 30\,\% and 37\,\% is indistinguishable from guessing.
\end{plainwords}

\section{How sure can you be about a score from 55 patches?}

\subsection{The Wilson interval}

An accuracy measured on 55 patches is an estimate of the accuracy the model would have on many more patches. The \textbf{Wilson score interval} \citep{wilson1927} gives a range that contains the true accuracy with 95\,\% confidence. With $\hat p = k/n$ correct answers and $z = 1.96$:
\begin{equation}
\frac{\hat p + \frac{z^2}{2n}}{1+\frac{z^2}{n}} \;\pm\; \frac{z}{1+\frac{z^2}{n}}\sqrt{\frac{\hat p(1-\hat p)}{n} + \frac{z^2}{4n^2}} .
\label{eq:wilson}
\end{equation}

\begin{table}[H]
\centering
\caption{Wilson 95\,\% intervals for the main accuracies on the 55 test patches (recomputed for this guide).}
\label{tab:wilson}
\begin{tabular}{lrr}
\toprule
Run & Accuracy & 95\,\% interval \\
\midrule
Qwen2-VL-2B zero-shot, full card & 56.4\,\% (31/55) & 43.3--68.6\,\% \\
Qwen2-VL-2B zero-shot, label line removed & 32.7\,\% (18/55) & 21.8--45.9\,\% \\
LLaVA-0.5B zero-shot, \texttt{compact} table & 43.6\,\% (24/55) & 31.4--56.7\,\% \\
Qwen2-VL-2B class LoRA, balanced, label-free & 36.4\,\% (20/55) & 24.9--49.6\,\% \\
Qwen2-VL-2B class LoRA, first run (superseded) & 67.3\,\% (37/55) & 54.1--78.2\,\% \\
Qwen2-VL-2B explanation LoRA & 61.8\,\% (34/55) & 48.6--73.5\,\% \\
LLaVA, BLIP-2 zero-shot & 32.7\,\% (18/55) & 21.8--45.9\,\% \\
Majority class & 34.5\,\% (19/55) & 23.4--47.7\,\% \\
$\Delta$NDVI rule fitted on training data & 94.5\,\% (52/55) & 85.1--98.1\,\% \\
\bottomrule
\end{tabular}
\end{table}

The intervals are about $\pm13$ percentage points wide. Differences of a few points between models mean nothing at this sample size.

\subsection{Comparing two models on the same patches: McNemar's test}

Because every model sees the same 55 patches, two models should be compared \emph{patch by patch}. \textbf{McNemar's test} \citep{mcnemar1947} uses only the patches where they disagree: $b$ patches only model~1 got right, $c$ patches only model~2 got right. If the two were equally good, $b$ would follow a binomial distribution with $b+c$ trials and probability 0.5, and the exact two-sided p-value is
\begin{equation}
p = \min\!\Bigl(1,\; 2\sum_{i=0}^{\min(b,c)}\binom{b+c}{i}\,0.5^{\,b+c}\Bigr).
\label{eq:mcnemar}
\end{equation}

\begin{table}[H]
\centering
\caption{Paired comparisons on the 55 test patches (recomputed for this guide).}
\label{tab:mcnemar}
\begin{tabular}{lrrr}
\toprule
Comparison & $b$ & $c$ & $p$ \\
\midrule
Qwen zero-shot (full card) vs.\ always ``increase'' & 31 & 19 & 0.12 \\
Qwen full card vs.\ Qwen without the label line (\texttt{compact}) & 18 & 5 & 0.011 \\
LLaVA \texttt{compact} vs.\ always ``increase'' & 16 & 11 & 0.44 \\
Balanced class LoRA vs.\ always ``increase'' & 20 & 19 & 1.0 \\
First class LoRA (superseded) vs.\ Qwen zero-shot & 19 & 13 & 0.38 \\
\bottomrule
\end{tabular}
\end{table}

\begin{plainwords}
Even Qwen's 56.4\,\% with the full card is not significantly better than always answering ``increase'' ($p = 0.12$), because Qwen never answered ``increase''. The one clearly significant difference is between Qwen with and without the card's summary line ($p = 0.011$). Without the line, neither LLaVA nor the balanced LoRA is distinguishable from the baseline. These facts are honest and worth a sentence each in the report.
\end{plainwords}

\section{Calibration (planned, not done)}

If a model states a confidence $q_i$ and is right ($o_i = 1$) or wrong ($o_i = 0$), its calibration can be scored with the Brier score $\frac{1}{N}\sum_i (q_i - o_i)^2$ or the expected calibration error \citep{brier1950,guo2017}. The roadmap planned this, but no model gave usable confidences: Qwen returned \texttt{"confidence": 0.0} in its free-form explanations, and the structured prompt's template value 0.5 was often copied unchanged. Calibration belongs in future work.

\section{The vocabulary of explanations}

\begin{longtable}{L{3.3cm}L{11.0cm}}
\caption{Terms from the roadmap (Section~3), with examples from this project.}\label{tab:xaiterms}\\
\toprule
\textbf{Term} & \textbf{Meaning and example} \\
\midrule
\endfirsthead
\toprule
\textbf{Term} & \textbf{Meaning and example} \\
\midrule
\endhead
Explanation & A human-readable account of the evidence for a conclusion, such as ``NDVI dropped from 0.625 to 0.575, indicating a decrease in green cover''. \\
Local and global & Local: why \emph{this} patch got \emph{this} answer. Global: how the model behaves over the whole test set (``always says decline''). \\
Post-hoc & Produced alongside or after a prediction. Every VLM justification here is post-hoc. \\
Plausibility & Whether an explanation \emph{sounds} reasonable. LLaVA's free-form text sounds plausible but contradicts the numbers. \\
Groundedness & Whether the claims are supported by the input. Checkable here, because the true changes are known. \\
Faithfulness & Whether the explanation reflects the process that actually produced the answer. Very hard to check for a black-box model. \\
Hallucination & A claim not supported by the input, such as ``rainfall increased'' when it fell. \\
Causal restraint & Not claiming causes the evidence cannot establish (``declined because of drought''). \\
\bottomrule
\end{longtable}

\subsection{Faithfulness is not groundedness}

\citet{jacovi2020} argue that a faithful explanation must reflect the model's real reasoning, and that a plausible text is not evidence of that. \citet{turpin2023} showed that language models write convincing step-by-step explanations that hide the real reason for an answer. Hallucination is a known weakness of VLMs: \citet{li2023pope} found that they often claim objects that are not in the image; here the equivalent is claiming changes that are not in the data.

\begin{plainwords}
When a model writes ``I chose decline because NDVI fell'', you cannot know whether NDVI really drove its answer. You can only check whether the sentence agrees with the data. That is groundedness (or evidence consistency), not faithfulness.
\end{plainwords}

\begin{foryourreport}
The project's code calls its score \texttt{faithfulness} (\file{src/explain_metrics.py}), but it measures \emph{evidence consistency}: the share of directional claims that agree with the measured changes. In the report, call it ``claim correctness'' or ``evidence consistency'' and say in one sentence why it is not faithfulness, citing \citet{jacovi2020}. This one sentence shows the examiner that you understand the difference.
\end{foryourreport}

\section{How an explanation is scored}
\label{sec:scorer}

The scorer is in \file{src/explain_metrics.py} and \file{scripts/score_explanations.py}.

\subsection{Step 1: find the claims}

The justification (and, for JSON answers, the supporting evidence, conflicting evidence, visual change and possible driver fields) is searched for \emph{indicator} words, each paired with the nearest \emph{direction} word:
\begin{itemize}
  \item indicators: NDVI, EVI, NDMI, NBR, ``green cover / greenness / canopy'', VV, ``VH / backscatter'', ``rainfall / precipitation / drier / wetter'', ``soil moisture'', ``temperature / warmer / cooler'';
  \item directions: increase-type (increase, rise, higher, gain, \dots), decrease-type (decrease, drop, decline, lower, loss, \dots), stable-type (stable, unchanged, little change, similar, \dots);
  \item the first direction word within 40 characters after the indicator is used; failing that, the closest one within 30 characters before it.
\end{itemize}

\subsection{Step 2: check each claim against the data}

The measured change for each indicator is turned into a direction using a dead-band $\tau$, below which a change counts as stable:
\begin{equation}
\mathrm{dir}(\Delta) =
\begin{cases}
\text{stable} & |\Delta| < \tau\\
\text{increase} & \Delta \ge \tau\\
\text{decrease} & \Delta \le -\tau
\end{cases}
\label{eq:deadband}
\end{equation}

\begin{table}[H]
\centering
\caption{Dead-band thresholds (\texttt{DIRECTION\_THRESHOLDS} in \file{src/explain_metrics.py}).}
\label{tab:deadband}
\begin{tabular}{lr|lr}
\toprule
Indicator & $\tau$ & Indicator & $\tau$ \\
\midrule
NDVI, EVI, NDMI, NBR & 0.02 & Rainfall & 5\,mm \\
Green cover (colour image) & 0.05 & Soil moisture & 0.01 \\
VV, VH & 0.3\,dB & Temperature & 0.3\,\degC \\
\bottomrule
\end{tabular}
\end{table}

A claim is \textbf{correct} when its direction equals the measured direction.

\subsection{Step 3: summary scores per model}

\begin{description}
  \item[Checkable claims per explanation] The mean number of claims that could be checked.
  \item[Correct claims] $\dfrac{\#\text{correct claims}}{\#\text{checkable claims}}$, pooled over all patches. The main explanation score.
  \item[Explanations with at least one wrong claim] Called \texttt{hallucination\_rate} in the code. It is strict: one wrong claim among five marks the whole explanation.
  \item[Self-consistency] Whether the most frequent claimed direction matches the model's own class (decline with decrease, increase with increase, stable with stable). Ties count as having no dominant direction and are left out (the original code broke ties at random, so the score changed between runs; this was fixed on 6 October).
  \item[ROUGE-L] Word overlap between the justification and a reference, from the longest common subsequence of words \citep{lin2004rouge}; computed against the evidence table and against the dataset's own narrative. High overlap means copying.
  \item[Distinct-2 and near-duplicate share] How varied a model's explanations are: unique word pairs divided by all word pairs, and the share of explanations that are almost identical (ROUGE-L above 0.9) to another one from the same model.
\end{description}

\begin{watchout}
Limits of this scorer that you should state:
\begin{itemize}
  \item It does not understand negation: ``NDVI did not decrease'' is read as a ``decrease'' claim.
  \item The nearest direction word may belong to a different indicator in the same sentence.
  \item The dead-band marks some true statements as wrong: ``EVI $\Delta$ $-0.006$, a decrease'' is wrong for the scorer because $0.006 < 0.02$.
  \item It cannot catch an invented number that falls inside the dead-band (for example, writing $+0.002$ when the change was $-0.005$).
  \item Self-consistency also counts climate claims; because rainfall fell everywhere, mentioning rainfall adds a ``decrease'' vote.
  \item A correct direction is not a good explanation: a model can list true numbers and still choose the wrong class.
\end{itemize}
\end{watchout}

\begin{foryourreport}
In ``Evaluation Metrics'' (about 450 words):
\begin{itemize}
  \item Accuracy, macro-F1 and the majority baseline, with Equations~\ref{eq:metrics}--\ref{eq:macrof1}.
  \item Wilson intervals (Equation~\ref{eq:wilson}) and McNemar's test for paired comparisons, with \citet{wilson1927} and \citet{mcnemar1947}.
  \item The three steps of the claim scorer, the dead-band table, and the four explanation scores (checkable claims, correct claims, at least one wrong claim, self-consistency).
  \item That the score measures evidence consistency, not faithfulness \citep{jacovi2020}; that ROUGE-L is reported only as a copying measure \citep{lin2004rouge}.
\end{itemize}
\end{foryourreport}

\begin{selfcheck}
\begin{enumerate}
  \item Why can a model with good accuracy still have a poor macro-F1?
  \item What does a Wilson interval of 43.3--68.6\,\% for Qwen's 56.4\,\% tell you?
  \item In McNemar's test, which patches are used, and why is Qwen's 56.4\,\% not significantly better than 34.5\,\%?
  \item What is the difference between groundedness and faithfulness? Which one does the project measure?
  \item Give two ways the claim scorer can mark a true statement as wrong.
\end{enumerate}
\end{selfcheck}
